% ===================================================================
% CHAPTER 10: CONCLUSION
% ===================================================================
\chapter{Conclusion}
\label{ch:conclusion}

\section{Summary of Contributions}
\label{sec:ch10:summary}

% Ported from: IDETC 2025 (DETC2025-167604), Sec. 7, transcribed from the published PDF
The testbed of Chapter~\ref{ch:tools} overcame previous testbed
limitations through HSV color space representation and neural network
calibration, achieving $R^2$ accuracies of 0.96 for direction and 0.91 for
magnitude while maintaining precise formation control of the robot
cluster. On it, the vector-sum and vector-to-scalar primitives ran in
continuous motion, confirming behaviors previously demonstrated only in
simulation.

% Ported from: cwaight_systems_paper_with_revisions2.tex
Chapter~\ref{ch:first_order} presented a framework for multirobot navigation in vector fields containing critical points. A Jacobian-based approach uses instantaneous measurements from three or more robots to identify critical point locations without prior field knowledge, supporting attraction and orbital control laws built on these estimates. Simulation across six canonical field types confirmed convergence of these control primitives. On curved fields the sign of $\det\hat{\mathbf{J}}$ on the first cycle decided which type of critical point the cluster reached. In 459 hardware runs, the cluster converged within 0.05~m in 99\% of saddle and 91\% of vortex trials with three robots and in every vortex trial with four, with mean position errors of 0.034~m or less. Orbital control experiments maintained commanded radii within a quarter of the radius from 0.20~m upward with three robots in both fields, and below that held bounded orbits that did not track the commanded radius; at 0.10~m the vortex orbit settled onto a limit cycle sized by the robots' delay and lag. Simulations matched to each run reproduced the saddle's attraction error and about half of the orbit fluctuation.

% Ported from: Draft_13.tex
Chapter~\ref{ch:second_order} showed that six robots sampling at once can fit a local
quadratic model of a planar flow, which gives the velocity gradient as a
field over the formation. Two surrogate fields read from that fit, the
Jacobian determinant $D$ and the smaller strain eigenvalue $s_1$, mark
coherent structures of the flow as trenches. On each field we built an
adaptive navigation primitive that acquires and rides those trenches from
instantaneous measurements, with no integration horizon and no
pre-straddling. Without noise, both acquire and traverse the
double-gyre separatrix from every start and hold the downstream saddle. Measured
by the noise level at which retention falls to 50\%, the $D$ tracker
withstands more measurement and position noise than the $s_1$ tracker from
every start tested. The $s_1$ tracker is objective past its
one seed step, and it alone holds the same path under a rotating observer.
A cluster with an absolute heading reference, such as a compass or GNSS,
should therefore carry the $D$ tracker, and one whose flow estimate
carries its own rotation should carry the $s_1$ tracker.

% New text. Primitive library, moved from Chapter 8 on 2026-09-30.
Table~\ref{tab:ch10:library} collects every primitive in this
dissertation with its fit order, the smallest cluster it was run with,
the quantity it steers on, its stability result, and how it was verified.

\begin{table}[htbp]
\centering
\footnotesize
\caption{Primitive library. Stability entries are as stated in the
chapter that defines each primitive.}
\label{tab:ch10:library}
\setlength{\tabcolsep}{3pt}
\begin{tabular}{>{\raggedright\arraybackslash}p{2.0cm} c c >{\raggedright\arraybackslash}p{2.7cm} >{\raggedright\arraybackslash}p{3.4cm} >{\raggedright\arraybackslash}p{3.0cm}}
\hline
Primitive & Order & Robots & Steers on & Stability & Verification \\
\hline
Vector-sum & 0 & 3 & mean sensed vector & none stated & hardware, fixed and sinking vortex, 10 runs each \\
Vector-to-scalar & 0 & 3 & slope of the sensed magnitude & none stated & hardware, fixed vortex, 10 runs plus randomized starts \\
Attraction & 1 & 3 & $\hat{\mathbf{p}}^* = -\hat{\mathbf{J}}^{-1}\hat{\mathbf{h}}$ & exponential convergence to a stationary estimate, $\tau = 1/k$ & simulation, 11 fields; hardware, 239 trials \\
Orbital & 1 & 3 & $\hat{\mathbf{p}}^*$ with radial and tangential terms & radial error $\dot{e}_r = -k_r e_r$ & simulation, 11 fields; hardware, 220 runs \\
$D$ tracker & 2 & 6 & $\hat{D}$ and $\hat{\mathbf{H}}_{D,0}$ & transverse practical stability, $\limsup|n| \leq \delta/a_\perp$ & simulation, double gyre and Santa Barbara Channel \\
$s_1$ tracker & 2 & 6 & $\hat{s}_1$ and $\nabla\hat{s}_1$ & same transverse bound, frame equivariant past its seed step & simulation, double gyre and Santa Barbara Channel \\
\hline
\end{tabular}
\end{table}

\section{Final Thoughts}
\label{sec:ch10:final}

% New text
The primitives in this dissertation are admittedly simple, reactive, and
minimal. Each uses the smallest cluster its fit order allows and only the
measurements of the current instant. Even in this form, however, they
show that the order of a cooperative fit sets which features of a flow a
cluster can reach, from a heading, to a critical point, to a separatrix.
Ongoing and future work targets a) hardware trials of the second-order
trackers on the Decabot testbed, b) control that accounts for a flow that
advects the robots, and c) a third-order fit from ten robots.
